% Generated draft: requires mathematical and source review.
\documentclass[11pt]{article}
\usepackage[margin=25mm]{geometry}
\usepackage{fontspec}
\setmainfont{Latin Modern Roman}
\usepackage{amsmath,amssymb,mathtools,unicode-math}
\setmathfont{Latin Modern Math}
\usepackage{enumitem,needspace,xurl,hyperref}
\setlength{\parindent}{0pt}
\setlength{\parskip}{5pt}
\setlength{\emergencystretch}{3em}
\newcommand{\N}{\mathbb{N}}
\newcommand{\Z}{\mathbb{Z}}
\newcommand{\Q}{\mathbb{Q}}
\newcommand{\R}{\mathbb{R}}
\newcommand{\C}{\mathbb{C}}
\newcommand{\F}{\mathbb{F}}
\newcommand{\A}{\mathbb{A}}
\newcommand{\PP}{\mathbb{P}}
\newcommand{\E}{\mathbb{E}}
\newcommand{\eps}{\varepsilon}
\newcommand{\dd}{\,\mathrm{d}}
\begin{document}
\section*{UnsolvedMath ID: 20000752}
\textbf{Source problem code:} catalogue entry
\subsection{Definitions and mathematical statement}
The original asks whether a sufficiently strong quantitative bounded-gap statement can imply a distribution level beyond Bombieri--Vinogradov. It attributes the converse question to V. Blomer, R. C. Vaughan's warning that the hypothesis would have to encode extensive $L$-function information, and K. Soundararajan's suggestion of a Hardy--Littlewood tuple conjecture with error $O(N^{1/2+\epsilon})$. The epsilon is reconstructed from a corrupted source character; tuple size, shift uniformity and weights are not specified.

For an explicit endpoint distribution convention, put $\psi(x;q,a)=\sum_{n\le x,\ n\equiv a\bmod q}\Lambda(n)$. Write $\operatorname{EH}_{\mathrm{end}}(\theta)$ for: for every $A>0$,
\[
\sum_{q\le x^\theta}\max_{(a,q)=1}\left|\psi(x;q,a)-\frac{x}{\varphi(q)}\right|
\ll_A \frac{x}{(\log x)^A}.
\]
This endpoint convention does not include a maximum over smaller endpoints $y\le x$.

The aid's moment reduction can be represented using the exact coprime-centered discrepancy
\[
\Delta_q(a;x)=\psi(x;q,a)-\frac1{\varphi(q)}\sum_{\substack{n\le x\\(n,q)=1}}\Lambda(n).
\]
For an integer $r\ge1$, a dyadic Gaussian-scale hypothesis is
\[
\sum_{R<q\le2R}\sum_{(a,q)=1}|\Delta_q(a;x)|^{2r}
\ll_r x^r R^{2-r}(\log x)^{B_r},
\]
uniformly over the required dyadic modulus ranges, with fixed $B_r$. The aid reports that this type of uniform hypothesis yields $\operatorname{EH}_{\mathrm{end}}(\theta)$ for every $\theta<r/(r+1)$, after the standard centering correction. Indeed its corresponding dyadic $L^1$ scale is $x^{1/2}R^{(r+1)/(2r)}$ times a logarithmic factor; the strict exponent inequality supplies a power saving. The exact source moment convention is not provided, so these definitions explicitly state the reconstruction used here.

The aid contrasts this with applying the triangle inequality separately to uniform square-root-size errors for distinct prime tuples. In that route, it reports that the above moment scale is controlled only for
\[
R\ge x^{1-1/(2r)+o(1)}.
\]
For $r>1$ this threshold lies above the useful upper exponent $r/(r+1)$. Thus separate tuple error bounds, summed absolutely over shift lattices, do not supply the required aggregate cancellation by this method. This is a limitation of that converse route, not a theorem that no stronger tuple hypothesis can imply EH.

Finally, the aid reports a periodic weighted model with the correct total mass and one sharply prescribed bounded-gap correlation but without equidistribution modulo $5$. Its weights and exact correlation identity are not supplied. It is offered only to show that total mass plus a single bounded-gap statistic need not force general residue equidistribution, not as a counterexample for the actual primes.
\subsection{Short English statement}
Determine how strong bounded-gap or prime-tuple information must be to improve the distribution level of primes. The supplied conditional reduction gives a higher-moment ladder, while its error-summation barrier identifies aggregate cancellation missing from a naive use of sharp individual tuple estimates.
\subsection{Sources}
\begin{itemize}
\item Local UnsolvedMath registry, permanent ID $20000752$, source code catalogue entry; canonical \texttt{aim-analytic-number-theory-notes.json}, \texttt{notes[115]}.
\item Angel Kumchev's notes, December 2005 AIM workshop ``Gaps between primes,'' problem 17: \url{https://aimath.org/WWN/primegaps/WorkshopProblems.pdf}. The record is reconstructed and unverified. It preserves Blomer's question and the comments of Vaughan and Soundararajan; Elliott--Halberstam, Bombieri--Vinogradov and Hardy--Littlewood are named without full references. No PDF inspection is claimed here.
\item Attached research attempt 1: conditional result, recorded run status open and rigor conditional, novelty confidence low; registry status partially solved, classification PARTIAL-PROGRESS, suggested difficulty L4. This machine-generated aid is not peer reviewed and requires independent verification before citation.
\end{itemize}
\end{document}
